%% 中文摘要页
\begin{ChineseAbstract}
近几年大语言模型技术研究突飞猛进，在软件开发领域也逐步从最初的代码补全发展到现在以 AI 智能体的角色贯穿从需求分析到测试交付的全流程。但受制于现有算力资源，智能体的上下文不是无限大的，目前主流模型的上下文上限大多在 1MB 左右，使得智能体在软件开发时需要一边自己搜索需要的文件内容，一边组织上下文，但这一过程通常需要分多次进行，会消耗大量的时间和 Token（大语言模型处理文本时的基本计量单位）。而且，组织好的上下文因为空间有限，在长对话过程中，如果模型上下文达到上限，就会挤掉最早开始时的上下文，造成“知识遗忘”，后续再次用到这些知识时，还需要重新搜索和组织。
针对这个问题，本文提出一种面向 AI 辅助开发的图谱驱动上下文工程方法。该方法首先从现有项目代码入手，建立代码模块、文件、调用方法之间的结构依赖关系索引，并将索引信息维护到本地数据库中。然后将现有项目的功能需求拆分成能力知识单元，通过串联知识单元与代码模块之间的映射关系，形成知识图谱。这样智能体在执行任务时，只需要调用现有知识库中知识单元的语义匹配，就能拿到基于当前需求语义对应的代码结构最小生成子图，而不是自己盲目检索，从而降低上下文组织成本。同时，在需求迭代过程中，不断新增和更新现有数据库中存储的知识单元与代码结构索引关系，确保知识图谱长期迭代后的可用性和准确性。
以上述方法为基础，实现了 WPW（Web Project Workflow，Web 项目开发工作流）系统。为了验证所提方法的有效性，本文以美食推荐系统为实验项目，分别使用 WPW 和 OpenSpec（开源的需求驱动开发工作流）从 0 到 1 完成系统开发，并对两种开发过程中的 Token（大语言模型处理文本时的基本计量单位）消耗进行对比。在当前项目和实验范围内，使用 WPW 的整体 Token 消耗更低。由于这一结果还不能直接说明差异来自上下文搜索和组织，因此进一步进行关键词检索实验，对比使用知识图谱和不使用知识图谱时的 Token 消耗。实验结果显示，使用知识图谱后，上下文搜索和组织过程中的 Token 消耗有所减少，表明知识图谱能够降低上下文管理带来的额外开销。

\ChineseKeywords{大语言模型；上下文工程；AI 辅助软件开发；代码知识图谱；美食推荐系统}
\end{ChineseAbstract}

%% 英文摘要页
\begin{EnglishAbstract}
In recent years, research on large language models has advanced rapidly. In software development, their role has expanded from the early use of code completion to AI agents that participate in the entire process from requirements analysis to testing and delivery. However, constrained by available computing resources, an agent's context is not unlimited. The context windows of most mainstream models are currently around 1 MB, so an agent working on a software project must search for the required file contents while also organizing them into context. This process often has to be repeated across multiple steps, consuming substantial time and tokens, where a token is the basic unit used by a large language model to process text. Moreover, because the organized context is also limited in size, a long conversation may reach the model's context limit and push out information from the beginning of the interaction. This causes the agent to lose previously acquired project knowledge and forces it to search for and organize the same knowledge again when it is needed later.

To address this problem, this thesis proposes a graph-driven context engineering method for AI-assisted software development. The method begins with the source code of an existing project and builds an index of structural dependencies among code modules, files, and methods, together with their call relationships, storing the index in a local database. It then decomposes the functional requirements of the project into capability knowledge units and connects these units to code modules through mapping relationships to form a knowledge graph. During task execution, the agent can use semantic matching over the capability units in the existing knowledge base to obtain a minimal subgraph of the code structure relevant to the current requirement, rather than searching through the project blindly. This reduces the cost of organizing task context. As requirements evolve, the method continuously adds and updates the capability units and code-structure index stored in the database, so that the knowledge graph remains usable and accurate over time.

Based on this method, we implement WPW (Web Project Workflow). To evaluate the method, we use a food recommendation system as the experimental project and develop it from scratch with both WPW and OpenSpec, an open-source requirements-driven development workflow. We compare the token consumption of the two development processes. Under the project and experimental conditions examined, the overall token consumption of WPW is lower. This comparison alone cannot establish that the difference comes from context search and organization, so we conduct an additional keyword-based retrieval experiment that compares token consumption with and without a knowledge graph. The results show that using a knowledge graph reduces token consumption during context search and organization, indicating that the knowledge graph can reduce the additional overhead introduced by context management.



\EnglishKeywords{Large Language Models; Context Engineering; AI-Assisted Software Development; Code Knowledge Graph; Food Recommendation System}
\end{EnglishAbstract}
